\section{Technologien und Implementierung}
\label{sec:technologien}

\subsection{Technologiestack}

\begin{table}[H]
\centering
\begin{tabular}{lll}
\toprule
\textbf{Technologie} & \textbf{Version} & \textbf{Verwendung} \\
\midrule
Python            & 3.11            & Alle Komponenten (Dispatcher, Worker, Client) \\
gRPC / grpcio     & $\geq$ 1.80     & Kommunikationsprotokoll \\
Protocol Buffers  & proto3          & Nachrichtenserialisierung \\
Docker            & 24+             & Containerisierung aller Komponenten \\
Docker Compose    & 2.x (v3.9)      & Orchestrierung des Gesamtsystems \\
\bottomrule
\end{tabular}
\caption{Technologiestack TaskGrid+}
\end{table}

\textbf{Wichtiger Hinweis zur grpcio-Version}:
Der generierte Proto-Code (\path{proto/taskgrid_pb2_grpc.py}) prüft
\texttt{GRPC\_GENERATED\_VERSION = '1.80.0'} zur Laufzeit und
wirft eine \texttt{RuntimeError} bei einer niedrigeren installierten Version.
Daher ist in allen \texttt{requirements.txt}-Dateien \texttt{grpcio>=1.80}
eingetragen (nicht nur \texttt{grpcio}).

\subsection{Strukturiertes Logging}

Alle Dispatcher-Log-Einträge folgen dem in §13 der Aufgabenstellung
definierten Format.
Das zentrale Logging-Modul \texttt{src/common/logger.py}
stellt zwei Funktionen bereit:

\begin{itemize}
  \item \texttt{get\_logger(component)}: Gibt einen konfigurierten Logger zurück.
        Gibt Logs gleichzeitig auf stdout (ohne Timestamp) und in eine Datei
        (mit Timestamp, Pfad: \path{LOG_DIR/<component>.log}) aus.
  \item \texttt{log\_event(logger, level, event, request\_id, task\_id, **extra)}:
        Erzeugt einen strukturierten Log-Eintrag im Format §13.
\end{itemize}

\textbf{Log-Format (Beispiele)}:

\begin{lstlisting}[language={},caption={Strukturierte Log-Ausgaben}]
[dispatcher] request_id=abc123 task_id=42 event=POST_TASK_accepted
             status=QUEUED type=sum queue_size=1
[dispatcher] request_id=abc123 task_id=42 event=DISPATCH_sending
             status=DISPATCHED worker_id=worker-sum-a1b2c3
[dispatcher] request_id=abc123 task_id=42 event=RESULT_RETURN_stored
             status=COMPLETED duration_ms=37
[dispatcher] task_id=42 event=DISPATCH_TIMEOUT_fired
             worker_id=worker-sum-a1b2c3 retry_count=1
[dispatcher] task_id=42 event=TIMEOUT_retrying
             retry_count=1 max_retries=3
\end{lstlisting}
Pflichtbestandteil des Testprotokolls: alle Ereignisse werden
mit \texttt{task\_id} und \texttt{request\_id} geloggt.

\noindent \textbf{Implementierungsdetail}: Python-Logger propagieren Nachrichten
standardmäßig die Hierarchie aufwärts.
Da jeder Logger eigene Handler hat (\texttt{dispatcher.dispatch\_loop}
$\rightarrow$ \texttt{dispatcher} $\rightarrow$ root), würde dies zu
doppelten Log-Einträgen führen.
Dies wurde durch \texttt{logger.propagate = False} in \texttt{get\_logger()} behoben.

\subsection{Dispatcher -- gRPC-Schnittstelle}

\texttt{PostTask} ist die primäre Eingangsschnittstelle des Dispatchers.
Die Validierungslogik prüft drei Bedingungen:

\begin{enumerate}
  \item \texttt{task\_type} darf nicht leer sein.
  \item \texttt{task\_type} muss $\leq$ 32 Zeichen lang sein.
  \item \texttt{task\_payload} muss $\leq$ 1024 Zeichen lang sein.
\end{enumerate}

\noindent Bei Verstößen wird der gRPC-Statuscode \texttt{INVALID\_ARGUMENT} gesetzt.
Valide Tasks erhalten eine neue \texttt{task\_id}, werden im
\texttt{TaskStore} gespeichert und in die \texttt{TaskQueue} eingereiht.

\subsection{Worker-Lookup}

Der Dispatcher löst Worker-Adressen nie statisch auf.
Stattdessen ruft \path{NamensdienstClient.lookup_worker()} bei jedem
Dispatch-Vorgang den Namensdienst ab.
Ein Round-Robin-Selector wählt einen Worker aus der zurückgegebenen Liste.
Bei leerer Liste wartet der Dispatcher \texttt{DISPATCH\_NO\_WORKER\_RETRY\_SECONDS}
und re-enqueued den Task.

\subsection{Ergebnis-Empfang}

\texttt{ReturnResult} empfängt Ergebnisse von Workern.
Drei Sonderfälle werden behandelt:

\begin{enumerate}
  \item \textbf{Unbekannte task\_id}: \texttt{Ack(success=False)} + Warning-Log.
  \item \textbf{Bereits terminaler Task}: Task wird ignoriert (Idempotenz §5).
        Dies verhindert Probleme bei verspäteten Ergebnissen nach Timeout.
  \item \textbf{Ungültiger Zustandsübergang}: \texttt{InvalidTransitionError}
        wird gefangen, Fehler wird geloggt.
\end{enumerate}

\noindent Bei erfolgreichem Empfang wird der Timeout-Timer sofort abgebrochen
(\path{dispatch_loop.cancel_timeout(task_id)}), bevor der
Zustandsübergang ausgeführt wird.
Die \texttt{duration\_ms} wird als Differenz zwischen
\texttt{timestamp\_dispatched} und \texttt{timestamp\_completed} berechnet
und im Log ausgegeben.

\subsection{Ergebnis-Abfrage}

\texttt{GetResult} gibt den aktuellen Zustand und -- falls vorhanden --
das Ergebnis eines Tasks zurück.
Für unbekannte \texttt{task\_id} wird gRPC-Code \texttt{NOT\_FOUND} gesetzt.

\subsection{Monitoring-Schnittstelle}

Der Dispatcher stellt zwei parallele Monitoring-Schnittstellen bereit:

\begin{enumerate}
  \item \textbf{gRPC \texttt{GetStatus}}: Gibt kompakten JSON-String zurück.
  \item \textbf{HTTP GET /status} (Port 8080): Gibt formatiertes JSON zurück,
        abrufbar mit \texttt{curl http://localhost:8080/status}.
\end{enumerate}

Der HTTP-Server läuft als Daemon-Thread (\texttt{HttpStatusServer}) neben
dem gRPC-Server.
Der \texttt{StatusCollector} aggregiert folgende Felder:

\begin{table}[H]
\centering
\begin{tabular}{m{6.5cm} m{10cm}}
\toprule
\textbf{Feld} & \textbf{Bedeutung} \\
\midrule
\texttt{aktive\_worker}                       & Anzahl aktiver Worker (via LookupWorker) \\
\texttt{unterstützte\_tasktypen}              & Alle bekannten Task-Typen \\
\texttt{offene\_tasks}                        & Tasks im Zustand QUEUED \\
\texttt{laufende\_tasks}                      & Tasks in DISPATCHED/PROCESSING/RETRYING \\
\texttt{abgeschlossene\_tasks}                & Tasks in COMPLETED \\
\texttt{fehlgeschlagene\_tasks}               & Tasks in FAILED/TIMEOUT \\
\path{durchschnittliche_bearbeitungszeit_ms} & Ø Bearbeitungszeit aller COMPLETED-Tasks \\
\texttt{anzahl\_timeouts}                     & Summe aller retry\_counts \\
\texttt{anzahl\_retries}                      & Identisch mit anzahl\_timeouts \\
\bottomrule
\end{tabular}
\caption{Monitoring-Felder (GET /status)}
\end{table}

\subsection{Timeout und Retry}

Für jeden an einen Worker gesendeten Task startet der Dispatcher
einen \texttt{threading.Timer}.
Die Timeout-Dauer ist über \texttt{DISPATCH\_TIMEOUT\_SECONDS} konfigurierbar
(default: 30\,Sekunden).

Der Timer-Callback \texttt{\_on\_timeout()} führt folgende Schritte aus:

\begin{enumerate}
  \item Frischen Task-Zustand aus dem Store lesen.
  \item Falls bereits terminal (COMPLETED/FAILED): Timer ignorieren
        (verspätetes Feuern -- kann passieren wenn Ergebnis kurz vor
        Timeout eintrifft).
  \item Zustandsübergänge: DISPATCHED $\rightarrow$ TIMEOUT $\rightarrow$ RETRYING.
  \item \texttt{retry\_count} inkrementieren.
  \item Falls \texttt{retry\_count} $\geq$ \texttt{MAX\_RETRIES}:
        RETRYING $\rightarrow$ FAILED (Task wird nicht erneut versucht).
  \item Sonst: Task in Queue re-enqueuen (RETRYING $\rightarrow$ DISPATCHED
        im nächsten Dispatch-Zyklus).
\end{enumerate}

\begin{lstlisting}[language=Python,caption={Timeout-Behandlung (dispatch\_loop.py)}]
if fresh.retry_count >= self._max_retries:
    transition(fresh, TaskState.FAILED)
    log_event(logger, "error", "TIMEOUT_max_retries_exceeded",
              task_id=task_id,
              retry_count=fresh.retry_count,
              max_retries=self._max_retries)
else:
    self._queue.enqueue(fresh)
    log_event(logger, "info", "TIMEOUT_retrying",
              task_id=task_id,
              retry_count=fresh.retry_count)
\end{lstlisting}

Der Timer wird abgebrochen wenn:
\begin{itemize}
  \item Ein Ergebnis via \texttt{ReturnResult} eintrifft (Normalfall).
  \item Der Worker die Task explizit ablehnt (ACK ok=False).
\end{itemize}

\subsection{Containerisierung}
\label{sec:containerisierung}

\subsubsection*{Verzeichnisstruktur der Container-Quellen}

Alle Komponenten liegen unterhalb von \texttt{src/} im Projekt-Root.
Jede Komponente bringt ihr eigenes \texttt{Dockerfile} mit:

\begin{lstlisting}[language={},caption={Container-relevante Verzeichnisstruktur}]
VerteilteSysteme-TaskGrid-/
  Dockerfile                    <- Dispatcher
  docker-compose.yml
  requirements.txt              <- grpcio>=1.80 (alle Komponenten)
  proto/                        <- gemeinsame gRPC-Stubs (alle Komponenten)
    taskgrid_pb2.py
    taskgrid_pb2_grpc.py
  src/
    dispatcher/                 <- Dispatcher-Quelltext
    namensdienst/
      Dockerfile                <- Namensdienst
      server.py                 <- Einstiegspunkt
      nameservice.py            <- Servicer-Implementierung
    worker/
      Dockerfile                <- Worker
      entrypoint.sh             <- ENV-VAR-Mapping + WORKER_ID-Generierung
      worker.py
      task_handlers/            <- Handler pro Tasktyp
    client/
      Dockerfile                <- Client
      client.py
    common/                     <- gemeinsamer Logger
\end{lstlisting}

\subsubsection*{Allgemeine Dockerfile-Prinzipien}

Alle Dockerfiles folgen denselben Entwurfsprinzipien:

\begin{enumerate}
  \item \textbf{Basis-Image \texttt{python:3.11-slim}}: Minimales Linux-Image
        ohne Desktop-Pakete. Reduziert Image-Größe und Angriffsfläche.

  \item \textbf{Layer-Cache-Optimierung}: \texttt{requirements.txt} wird als
        erster Schritt kopiert und installiert, der Quelltext erst danach.
        Docker cached den \texttt{pip install}-Layer solange
        \texttt{requirements.txt} unverändert bleibt -- auch wenn sich
        Quelltextdateien ändern. Damit entfällt die ca. 30-sekündige
        grpcio-Installation bei jedem Rebuild.

  \item \textbf{PYTHONPATH=/app:/app/proto}: Alle Komponenten verwenden den
        Projekt-Root als Python-Basis. \texttt{/app/proto} wird zusätzlich
        benötigt, weil der generierte gRPC-Stub
        (\texttt{taskgrid\_pb2\_grpc.py}) seit grpcio\,$\geq$\,1.80 einen
        absoluten Import \texttt{import taskgrid\_pb2} (ohne Paket-Prefix)
        enthält. Ohne \texttt{/app/proto} im Suchpfad würde dieser Import
        zur Laufzeit mit \texttt{ModuleNotFoundError} fehlschlagen.

  \item \textbf{Getrennte Build-Kontexte}: Alle Container nutzen den
        Projekt-Root als Build-Kontext (\texttt{context: .}) und referenzieren
        ihr Dockerfile per \texttt{dockerfile:}-Pfad. Dadurch können alle
        Komponenten die gemeinsamen \texttt{proto/}-Stubs ohne Duplizierung
        einbinden.

  \item \textbf{Umgebungsvariablen als Dokumentation}: Jede konfigurierbare
        Größe erscheint mit ihrem Default im Dockerfile. Überschreibungen
        erfolgen ausschließlich in \texttt{docker-compose.yml}.

  \item \textbf{Log-Verzeichnis \texttt{/app/logs}}: Wird per \texttt{RUN mkdir}
        vorab angelegt und via Volume in docker-compose.yml gemountet
        (Aufgabenanforderung §8: Volumes ausschließlich für Logs).
\end{enumerate}

\subsubsection*{Dispatcher Dockerfile}

Build-Kontext: Projekt-Root. Benötigt \texttt{proto/} (gRPC-Stubs) und
\texttt{src/} (Dispatcher- und Common-Module).

\begin{lstlisting}[language={},caption={Dispatcher Dockerfile (src/: Dockerfile)}]
FROM python:3.11-slim
WORKDIR /app
COPY requirements.txt .
RUN pip install --no-cache-dir -r requirements.txt
COPY proto/ ./proto/
COPY src/   ./src/
ENV PYTHONPATH=/app:/app/proto
ENV DISPATCHER_PORT=50051
ENV STATUS_HTTP_PORT=8080
ENV NAMENSDIENST_HOST=namensdienst
ENV NAMENSDIENST_PORT=50052
ENV DISPATCH_TIMEOUT_SECONDS=30
ENV MAX_RETRIES=3
ENV DISPATCH_NO_WORKER_RETRY_SECONDS=2
EXPOSE 50051
EXPOSE 8080
CMD ["python", "src/dispatcher/server.py"]
\end{lstlisting}

\noindent \textbf{Keine Worker-Adressen}: Gemäß Aufgabenanforderung §4.2 enthält das
Dockerfile keine Worker-Adressen oder -Ports. Der Dispatcher ermittelt
Worker ausschließlich dynamisch über den Namensdienst (\texttt{LOOKUP\_WORKER}).
\newpage
\subsubsection*{Client Dockerfile}

Build-Kontext: Projekt-Root. Benötigt \texttt{proto/} und \texttt{src/client/}.

\begin{lstlisting}[language={},caption={Client Dockerfile (src/client/Dockerfile)}]
FROM python:3.11-slim
WORKDIR /app
COPY requirements.txt .
RUN pip install --no-cache-dir -r requirements.txt
COPY proto/       ./proto/
COPY src/client/  ./src/client/
ENV PYTHONPATH=/app:/app/proto
ENV DISPATCHER_ADDRESS=dispatcher:50051
ENV CLIENT_ID=client-1
EXPOSE 50051
CMD ["python", "src/client/client.py"]
\end{lstlisting}

\subsubsection*{Worker Dockerfile}

Build-Kontext: Projekt-Root. Benötigt \texttt{proto/} (gRPC-Stubs, geteilt
mit Dispatcher), \path{src/worker/task_handlers/} und
\texttt{src/worker/worker.py}.

\begin{lstlisting}[language={},caption={Worker Dockerfile (src/worker/Dockerfile)}]
FROM python:3.11-slim
WORKDIR /app
COPY src/worker/requirements.txt .
RUN pip install --no-cache-dir -r requirements.txt
COPY proto/                    ./proto/
COPY src/worker/task_handlers/ ./src/worker/task_handlers/
COPY src/worker/worker.py      ./src/worker/worker.py
COPY src/worker/entrypoint.sh  .
RUN sed -i 's/\r//' entrypoint.sh && chmod +x entrypoint.sh
ENV PYTHONPATH=/app:/app/proto
ENV WORKER_TYPE=reverse,sum,hash,upper,wait
ENV WORKER_HOST=worker
ENV WORKER_PORT=50051
ENV NAMENSDIENST_ADDRESS=namensdienst:50052
ENV DISPATCHER_ADDRESS=dispatcher:50051
ENV HEARTBEAT_INTERVAL_SECONDS=5
EXPOSE 50051
ENTRYPOINT ["./entrypoint.sh"]
\end{lstlisting}
\textbf{sed-Kommando}: Das Repository wird auf Windows entwickelt.
Windows-Zeilenenden (CRLF) in Shell-Skripten machen den Shebang
(\texttt{\#!/bin/sh}) im Linux-Container unlesbar, da \texttt{\textbackslash r}
als Teil des Interpreter-Pfads interpretiert wird.
\texttt{sed -i 's/\textbackslash r//'} entfernt alle \texttt{\textbackslash r}-Zeichen
bevor dem Skript Ausführungsrechte gegeben werden.

\noindent \textbf{entrypoint.sh -- ENV-VAR-Mapping und WORKER\_ID-Generierung}:
Die Issue-Spezifikation (\#36) und \texttt{worker.py} nutzen unterschiedliche
Variablennamen. Das Entrypoint-Skript übersetzt zur Laufzeit:

\begin{lstlisting}[language=bash,caption={src/worker/entrypoint.sh}]
#!/bin/sh
# Issue-Namen -> worker.py-interne Namen
export TASK_TYPES="${WORKER_TYPE:-${TASK_TYPES:-reverse}}"
export NAMING_SERVICE_ADDRESS="${NAMENSDIENST_ADDRESS:-namensdienst:50052}"
export HEARTBEAT_INTERVAL="${HEARTBEAT_INTERVAL_SECONDS:-5}"

# Eindeutige ID pro Container-Instanz (Docker setzt hostname = Container-ID-Prefix)
if [ -z "${WORKER_ID}" ]; then
    export WORKER_ID="worker-${TASK_TYPES}-$(hostname)"
fi

exec python src/worker/worker.py
\end{lstlisting}
Die \texttt{WORKER\_ID}-Generierung via \texttt{hostname} garantiert, dass
bei \texttt{docker compose up --scale worker-sum=3} alle drei Instanzen
unterschiedliche IDs beim Namensdienst registrieren.
\texttt{exec} ersetzt den Shell-Prozess durch Python, sodass
SIGTERM-Signale (Docker-Stop) direkt beim Worker ankommen.

\subsubsection*{Namensdienst Dockerfile}

Build-Kontext: Projekt-Root. Benötigt \texttt{proto/},
\texttt{src/namensdienst/} und \texttt{src/common/} (Logger).

\begin{lstlisting}[language={},caption={Namensdienst Dockerfile (src/namensdienst/Dockerfile)}]
FROM python:3.11-slim
WORKDIR /app
COPY src/namensdienst/requirements.txt .
RUN pip install --no-cache-dir -r requirements.txt
COPY proto/             ./proto/
COPY src/namensdienst/  ./src/namensdienst/
COPY src/common/        ./src/common/
ENV PYTHONPATH=/app:/app/proto
ENV NAMESERVICE_PORT=50052
ENV HEARTBEAT_TIMEOUT_SECONDS=15
ENV HEARTBEAT_CHECK_INTERVAL_SECONDS=5
EXPOSE 50052
CMD ["python", "src/namensdienst/server.py"]
\end{lstlisting}

Der Einstiegspunkt ist \texttt{server.py}, der den gRPC-Server startet und
die \texttt{Namensdienst}-Klasse aus \texttt{nameservice.py} als Servicer
registriert.

\subsubsection*{docker-compose.yml}

Die Orchestrierungsdatei steuert alle acht Container-Services:

\begin{table}[H]
\centering
\begin{tabular}{llll}
\toprule
\textbf{Service} & \textbf{Dockerfile} & \textbf{Port (extern)} & \textbf{Zweck} \\
\midrule
\texttt{namensdienst} & \texttt{src/namensdienst/Dockerfile} & 50052 & Service Discovery \\
\texttt{dispatcher}   & \texttt{Dockerfile}                  & 50051, 8080 & Task-Verwaltung \\
\texttt{worker-reverse} & \texttt{src/worker/Dockerfile}     & -- & Tasktyp reverse \\
\texttt{worker-sum}     & \texttt{src/worker/Dockerfile}     & -- & Tasktyp sum \\
\texttt{worker-hash}    & \texttt{src/worker/Dockerfile}     & -- & Tasktyp hash \\
\texttt{worker-upper}   & \texttt{src/worker/Dockerfile}     & -- & Tasktyp upper \\
\texttt{worker-wait}    & \texttt{src/worker/Dockerfile}     & -- & Tasktyp wait \\
\texttt{client}         & \texttt{src/client/Dockerfile}     & -- & One-Shot Demo \\
\bottomrule
\end{tabular}
\caption{docker-compose.yml -- Services im Überblick}
\end{table}
\newpage
\textbf{Implementierte Anforderungen}:

\begin{table}[H]
\centering
\begin{tabular}{p{4.5cm}p{10cm}l}
\toprule
\textbf{Anforderung} & \textbf{Umsetzung} & \textbf{Ref.} \\
\midrule
Startreihenfolge &
  \texttt{depends\_on: condition: service\_healthy} -- Worker starten erst
  wenn Namensdienst \emph{und} Dispatcher healthy sind &
  \#38 \\
Health Checks &
  TCP-Port-Sonde als Python-Einzeiler (kein \texttt{netcat}/\texttt{curl}
  in slim-Image nötig) &
  \#38 \\
Typspezifische Worker &
  Separate Services \texttt{worker-sum}, \texttt{worker-reverse} etc.,
  jeweils mit \texttt{WORKER\_TYPE}-ENV &
  \#38 \\
Skalierung &
  Kein \texttt{container\_name} bei Workers $\Rightarrow$
  \texttt{docker compose up --scale worker-sum=3} möglich &
  \#39 \\
Volumes &
  Ausschließlich \texttt{logs}-Volume (kein gemeinsames Volume
  für Task-/Ergebnisdaten) &
  §8 \\
Internes Netzwerk &
  Bridge-Netzwerk \texttt{taskgrid}; Worker-Ports nicht nach außen
  exponiert &
  §11 \\
One-Shot Client &
  \texttt{restart: "no"} -- Client beendet sich nach Demo &
  §11 \\
Konfiguration &
  Ausschließlich Umgebungsvariablen, keine hartcodierten Adressen &
  §4.2 \\
\bottomrule
\end{tabular}
\caption{docker-compose.yml -- Erfüllte Anforderungen}
\end{table}

\textbf{Startreihenfolge}:

\begin{lstlisting}[language={},caption={Startabhängigkeiten in docker-compose.yml}]
namensdienst  --[healthy]--> dispatcher
                             dispatcher  --[healthy]--> worker-*
                                                        worker-*  -->  (client)
\end{lstlisting}

\textbf{Health Check -- Implementierungsdetail}:
\texttt{python:3.11-slim} enthält kein \texttt{netcat} oder \texttt{curl}.
Stattdessen wird ein Python-Einzeiler als Health Check eingesetzt:

\begin{lstlisting}[language=yaml,caption={Health Check (Dispatcher, Port 50051)}]
healthcheck:
  test: ["CMD", "python", "-c",
         "import socket; s=socket.socket(); s.settimeout(2);
          s.connect(('localhost', 50051)); s.close()"]
  interval: 5s
  timeout: 3s
  retries: 12
  start_period: 10s
\end{lstlisting}

Dieser Check prüft, ob der gRPC-Port eine TCP-Verbindung akzeptiert --
nicht nur ob der Python-Prozess gestartet wurde.
\texttt{start\_period: 10s} räumt dem Dispatcher Zeit ein, um seinen
eigenen Startup (Verbindung zum Namensdienst, Thread-Start) abzuschließen,
bevor fehlgeschlagene Checks als Fehler zählen.

\textbf{Skalierung}:

Einzelne Worker-Typen können horizontal skaliert werden ohne
Konfigurationsänderungen:

\begin{lstlisting}[language=bash,caption={Skalierungsbeispiele}]
# Drei parallele sum-Worker starten
docker compose up --scale worker-sum=3

# Mehrere Typen gleichzeitig skalieren
docker compose up --scale worker-reverse=2 --scale worker-hash=2

# Gesamtsystem starten (je ein Worker pro Typ)
docker compose up
\end{lstlisting}

Die Skalierung funktioniert, weil:
\begin{itemize}
  \item Worker-Services haben kein \texttt{container\_name} -- Docker vergibt
        automatisch eindeutige Namen.
  \item \texttt{entrypoint.sh} generiert \texttt{WORKER\_ID} via
        \texttt{hostname} (= Container-ID-Prefix) -- keine Namenskollisionen
        bei mehreren Instanzen desselben Typs.
  \item Docker DNS löst den Service-Namen (z.\,B. \texttt{worker-sum}) per
        Round-Robin auf alle laufenden Instanzen auf. Alle registrieren sich
        beim Namensdienst mit demselben \texttt{WORKER\_HOST}, der Dispatcher
        erhält sie über \texttt{LOOKUP\_WORKER} zurück.
\end{itemize}
\newpage
\textbf{Kommunikationsports im Netzwerk \texttt{taskgrid}}:

\begin{table}[H]
\centering
\begin{tabular}{lll}
\toprule
\textbf{Verbindung} & \textbf{Port} & \textbf{Protokoll} \\
\midrule
Client $\rightarrow$ Dispatcher & 50051 & gRPC \\
Dispatcher $\rightarrow$ Namensdienst & 50052 & gRPC \\
Worker $\rightarrow$ Namensdienst & 50052 & gRPC \\
Dispatcher $\rightarrow$ Worker & 50051 & gRPC \\
Worker $\rightarrow$ Dispatcher & 50051 & gRPC \\
Monitoring $\rightarrow$ Dispatcher & 8080 & HTTP \\
\bottomrule
\end{tabular}
\caption{Interne Kommunikationsports im Docker-Netzwerk}
\end{table}
